\section{Módulos libres, generación y propiedades universales}\label{sec:modulos-libres}

\subsection{Construcción de módulos libres y generación}\label{tema:14-3-1}\label{tema:15-1-1a}
\begin{defi}
Un \emph{módulo libre sobre un conjunto $A$} es un módulo $F^R(A)$ con una función $j:A\to F^R(A)$ tal que toda función $f:A\to M$ se extiende de manera única a un morfismo de módulos $\varphi:F^R(A)\to M$, con $\varphi j=f$.
\end{defi}
Si existen dos objetos así, los mediadores correspondientes a sus funciones de generadores dan morfismos en ambos sentidos. Sus composiciones fijan los generadores; por unicidad son las identidades. Por ello el objeto libre es único salvo el isomorfismo que respeta esos generadores.

Para un módulo $N$ y un conjunto $A$, definimos
\[
N^{\oplus A}=\set{u:A\to N\given\supp u\text{ es finito}},
\qquad \supp u=\set{a\in A\given u(a)\ne0}.
\]
Con suma y escalares puntuales es módulo: el soporte de una suma está contenido en la unión de los soportes y el de $ru$ está contenido en el de $u$.

\begin{prop}\label{prop:C14-libre}
El módulo $R^{\oplus A}$ es libre sobre $A$, con $j(a)=e_a$, donde $e_a(x)=1_R$ si $x=a$ y $e_a(x)=0$ si $x\ne a$.
\end{prop}
\begin{proof}
Toda $u\in R^{\oplus A}$ tiene una expresión finita
\[
u=\sum_{a\in A}u(a)e_a.
\]
Una extensión lineal de $f:A\to M$ está forzada a satisfacer
\[
\varphi(u)=\sum_{a\in A}u(a)f(a).
\]
La fórmula está definida porque el soporte es finito. Preserva sumas por distributividad y, para $r\in R$, $\varphi(ru)=\sum(ru(a))f(a)=r\sum u(a)f(a)$. También $\varphi(e_a)=f(a)$. Esto prueba existencia y unicidad, sin suponer conmutatividad de $R$:
\[
\begin{tikzcd}[column sep=large,row sep=large]
R^{\oplus A}\arrow[r,dashed,destacado,"\exists!\,\varphi"]&M\\
A\arrow[u,"j"]\arrow[ur,"f"']&
\end{tikzcd}
\qedhere
\]
\end{proof}
Si $A=\{1,\ldots,n\}$, se obtiene $R^{\oplus A}=R^n$ y $e_i$ tiene un uno en la coordenada $i$ y ceros en las demás. Para $A$ vacío, el módulo libre es cero.

\begin{ejer}[Inyectividad de los generadores libres]\label{ejer:C14-generadores}
Suponiendo $R\ne0$, comprobar que la función $j$ de un módulo libre es inyectiva. Examinar también el caso del anillo cero.
\end{ejer}
\begin{sol}
En el modelo construido, si $a\ne b$, entonces $e_a(a)=1_R\ne0=e_b(a)$, luego $e_a\ne e_b$. La unicidad del objeto libre transporta la inyectividad a cualquier otro modelo. También se deduce directamente de la propiedad universal eligiendo $f:A\to R$ con $f(a)=1_R$, $f(b)=0$: si $j(a)=j(b)$, toda extensión daría $f(a)=f(b)$, contradicción.

Si $R=0$, cualquier módulo unitario es cero, porque $m=1_Rm=0_Rm=0$. El módulo libre sigue existiendo y es cero para cualquier $A$. En ese caso, $j$ es inyectiva exactamente cuando $|A|\leq1$.
\end{sol}
\begin{obs}
La construcción libre también vale para $R=0$. En ese caso $R^{\oplus A}=0$, y la función de generadores $j:A\to0$ sólo es inyectiva cuando $|A|\leq1$.
\end{obs}

Si $A\subseteq M$, la inclusión de conjuntos se extiende de manera única a un morfismo $\varphi_A:R^{\oplus A}\to M$. Su imagen es el \emph{submódulo generado por $A$}, denotado $\ang{A}$:
\[
\ang{A}=\set*{\sum_{a\in A}r_aa\given
 r_a\in R,\ r_a=0\text{ salvo para finitos }a}.
\]
Contiene $A$, porque $a=1_Ra$, y todo submódulo que contiene $A$ contiene todas esas combinaciones. Es, por tanto, el menor submódulo con esa propiedad. Para $A=\varnothing$ se obtiene cero.

\begin{defi}
Un módulo $M$ es \emph{finitamente generado} si $M=\ang{a_1,\ldots,a_n}$ para alguna familia finita. Equivalentemente, existe un homomorfismo sobreyectivo $R^n\to M$ para algún $n\geq0$.
\end{defi}
La equivalencia se sigue de la fórmula del mediador libre: $(r_1,\ldots,r_n)\mapsto\sum_i r_ia_i$. Recíprocamente, las imágenes de los vectores coordenados generan la imagen de todo morfismo desde $R^n$.
